\section*{Change Log}
\phantomsection
\addcontentsline{toc}{section}{Change Log}
\label{sec:change-log}

This page records the analysis-level changes between circulated versions of the note.
It is not a substitute for the configuration and figure provenance given in the body
and appendices. Numerical results from different versions should not be combined: each
version corresponds to a distinct, internally consistent analysis state.

\begingroup
\footnotesize
\setlength{\tabcolsep}{4pt}
\renewcommand{\arraystretch}{1.04}
\setlength{\LTleft}{\fill}
\setlength{\LTright}{\fill}
\begin{longtable}{P{0.11\textwidth}P{0.16\textwidth}P{0.64\textwidth}}
\caption{Analysis-note version history. The v4.2 row defines the scope of the
present document. Earlier versions are summarized to keep changes in inputs,
validation, and statistical interpretation traceable to reviewers.}
\label{tab:note-change-log}\\
\toprule
Version & Circulation date & Principal changes \\
\midrule
\endfirsthead
\multicolumn{3}{l}{\footnotesize\itshape Table~\ref{tab:note-change-log}
continued from the previous page.}\\
\toprule
Version & Circulation date & Principal changes \\
\midrule
\endhead
\midrule
\multicolumn{3}{r}{\footnotesize\itshape Continued on the next page.}\\
\endfoot
\bottomrule
\endlastfoot
v1 & March--May 2026 &
Established the full-range Gaussian-process background model, the blinded local
signal-extraction likelihood, the radiative-fraction conversion to $\eps^2$, and the
first single-dataset and combined validation studies. The figures in this version were
development results and did not define the present candidate configuration. \\
\addlinespace
v2 & June 2026 &
Introduced the $2.25\,\sigma_m$ blind/training geometry, matched-reference full-refit
pseudoexperiments, the shared-$\eps^2$ simultaneous likelihood, and expanded 2015,
2016, and 2021 method documentation. This version used common or development-stage
length-scale choices and retained broader engineering-run scan ranges in several
studies. \\
\addlinespace
v3 & July 10, 2026 &
Uses the regenerated 2021 1\% high-$p_{\mathrm{sum}}$ histogram, with no vertex-fit
$\chi^2$ requirement and at least eight electron-track hits; restricts the 2015 and
2016 scans to \SIrange{19}{90}{MeV} and \SIrange{39}{180}{MeV}; freezes the
dataset-dependent lower length-scale factors at 1.0, 0.9, and 1.1; and adopts 90\% CL
upper limits obtained with \CLs{}. The main validation section now contains the
functional-form seed construction and the injected-signal lower-bound scans, while
superseded closure and limit studies are retained with their original context in the
appendices. The results section documents the individual and simultaneous limits,
local and approximate global $p$-values, and all remaining pre-unblinding audit
qualifications. \\
\addlinespace
v4 & August 3, 2026 &
Separates the GP fit support from the signal-search interval in all three campaigns.
The 2015 and 2016 searches keep their reviewed mass grids but regain sidebands beyond
both endpoints, while the 2021 10\% sample retains the wider support developed in the
recent endpoint study. The result now uses the full 2015 and 2016 datasets in the
shared-$\eps^2$ fit and includes conditional background-only expected bands. The same
pseudoexperiments place the observed limit within the ensemble through the strong-,
weak-, and bounded two-sided consistency tests used in v3. Those quantities remain
upper-limit diagnostics rather than discovery significances or a substitute for
full-procedure coverage. \\
\addlinespace
v4.1 & August 4--5, 2026 &
Responds to the full-2016 length-scale saturation with a controlled observed-only
scan of $k_{\max,2016}=8,10,12,15,$ and 20. Factor 12 is the first nonbinding value
and is followed by a stable plateau. The exact combined observed limit and local
asymptotic $p_0$ scan are recomputed with all other analysis settings fixed. A
ten-toy 2021 source-and-exposure screen---native 1\%, $1\%\times10$,
$1\%\times100$, native 10\%, and $10\%\times10$---is included only as a fit-space
hyperparameter pilot. Its reviewed background-only scan supports factor 20 for the two
projected-100\% lanes and factor 25 if one bound must cover every pilot exposure;
the fixed-amplitude screen shows no factor-20 sensitivity loss, while a common 2--3\%
under-response remains an unresolved ten-toy closure diagnostic. No new expected
bands, limit-tail ensemble result, toy-calibrated discovery claim, or coverage claim
is included; the observed result remains a candidate pending production-matched
closure and direct coverage. The August 5 methodology revision adds explanatory
schematics for $C$, $\ell$, fixed $\alpha_i$, and the per-hypothesis optimization;
it changes no configuration or numerical result. \\
\addlinespace
v4.2 & August 5, 2026 &
Accepts $k_{\max,2016}=12$ as the current full-2016 setting and completes the
232-mass combined 90\% \CLs{} search with exactly 300 conditional background-only
joint pseudoexperiments at each mass. Central 68\% and 95\% intervals are
reported for that principal shared-$\eps^2$ search. The note also adds
standalone observed limits for the individual 2015 full, 2016 full, and 2021
10\% distributions, auxiliary 100-pseudoexperiment standalone and pairwise
conditional limit bands, the combined strong-, weak-, and bounded two-sided
limit-tail diagnostics, and combined and constituent local asymptotic $p_0$
plots. The auxiliary ensembles select indices 0--99 from the accepted
300-draw parent stream and reuse campaign pseudo-spectra across scopes. No
scan-toy-calibrated global discovery probability or direct coverage claim is
introduced; all bands remain descriptive quantiles conditional on the reviewed
fixed GP states. The results section also
adds an exact fixed-state 65~MeV extraction, including sideband context,
background-subtracted spectra, conditional Pearson-residual diagnostics, and a
comparison of the three standalone signal estimates with the simultaneous
shared-$\eps^2$ fit. \\
\end{longtable}
\endgroup
